\begin{example}\label{rl-0012}
  Suppose \((j_i:\mathfrak C\to \mathfrak E_i)\) is an embedding of the intrinsic normal cone stack \(\mathfrak C\) for \(\eta:Z\to X\) into a vector bundle stack \(\mathfrak E_i\) where \(i = 1,2\). Then it is not necessarily the case that the diagonal map \(j=(j_1,j_2):\mathfrak C\to \mathfrak E_1\oplus \mathfrak E_2\) is a closed embedding.

  To see this, let \(Z = X = \Spec k\), \(V = \mathbb A^1\) be the rank 1 vector bundle over \(Z\) and \(\mathfrak C = \mathfrak E_1 = \mathfrak E_2 = [0 / V] = BV\). Note that \(\mathfrak C\) is the intrinsic normal cone of the morphism \(\Spec k \to [\Spec k/V]\). Take \(j_i:\mathfrak C\to \mathfrak E_i\) to be the identity for both \(i\). The diagonal is then
  \[j = (j_1,j_2): BV \to BV \oplus BV = B(V\oplus V)\]
  induced by the diagonal map \(V\to V\oplus V\). The fiber of a point \(\Spec k\) over this diagonal is
  \[\Spec k\times_{B(V\oplus V)}BV = V,\]
  so the map \(j\) has positive dimensional fibers.

  On the level of complexes, this can be seen as a failure of \(h^0\) to be an isomorphism (as per the criterion in \cite[Proposition 2.6]{behrend-fantechi_IntrinsicNormalCone1997}). The two-term truncation of the cotangent complex of \(f:\Spec k \to BV\) is \(L_f = [0 \to k]\) (in degree \(-1\) and \(0\)) and the complex corresponding to \(B(V\oplus V)\) is \([0 \to k^2]\) (dualize and quotient to get \(B(V\oplus V)\)). The candidate obstruction theory is then determined by the morphism in degree \(0\), which is \((a,b) \mapsto a+b\), and this is exactly the map \(h^0(\phi)\) as well. This is surjective, but not an isomorphism, hence \(\phi^\vee:h^1/h^0(L_f^\vee)\to h^1/h^0\big((E_1^\bullet\oplus E_2^\bullet)^\vee\big)\) fails to be a closed embedding.
\end{example}
